\section{Negación con sustantivos}
\textbf{Indefinidos y cantidad cero.} Usa \textit{geen} para negar existencia o posesión cuando no hay artículo definido.

\begin{tabular}{@{}lll@{}}
\toprule
Afirmación & Negación & Traducción \\
\midrule
\textit{Hij heeft een fiets} & \textit{Hij heeft geen fiets} & Él no tiene bici \\
\textit{We hebben melk} & \textit{We hebben geen melk} & No tenemos leche \\
\textit{Er zijn stoelen} & \textit{Er zijn geen stoelen} & No hay sillas \\
\bottomrule
\end{tabular}

\textbf{Sustantivos definidos o posesivos.} No uses \textit{geen}. Niega con \textit{niet} porque hablas de algo concreto: \textit{Het is niet mijn boek}. 

\textbf{Mini práctica}
\begin{itemize}[leftmargin=*]
  \item Cambia a negación: \textit{Zij heeft een kat}. 
  \item Niega con posesivo: \textit{Dat is mijn fiets}. (Mantén \textit{niet}).
\end{itemize}
